\section{Signs of permutations}

The notion of the \emph{sign} (aka \emph{signature}) of a permutation is a
simple consequence of that of its length; moreover, it is rather well-known,
due to its role in the definition of a determinant. Thus we will survey its
properties quickly and without proofs.

\subsection{Definition and basic properties}

\begin{definition}
\label{def.perm.sign}
\lean{Equiv.Perm.sign_eq_neg_one_pow_invCount, Equiv.Perm.sign_coe_eq_neg_one_pow_invCount}
\leantarget
\leanok
Let $n \in \mathbb{N}$. The \emph{sign} of a permutation
$\sigma \in S_n$ is defined to be the integer $(-1)^{\ell(\sigma)}$.

It is denoted by $(-1)^{\sigma}$ or $\operatorname{sgn}(\sigma)$ or
$\operatorname{sign}(\sigma)$ or $\varepsilon(\sigma)$. It is also known
as the \emph{signature} of $\sigma$.
\end{definition}

\begin{proposition}
\label{prop.perm.sign.props}
\lean{Equiv.Perm.sign_id, Equiv.Perm.sign_transposition, Equiv.Perm.sign_isCycle, Equiv.Perm.sign_mul', Equiv.Perm.sign_prod_list, Equiv.Perm.sign_inv', Equiv.Perm.sign_eq_prod_pairs, Equiv.Perm.prod_diff_comp_perm}
\leantarget
\leanok
Let $n \in \mathbb{N}$.

\textbf{(a)} The sign of the identity permutation $\operatorname{id} \in S_n$
is $(-1)^{\operatorname{id}} = 1$.

\textbf{(b)} For any two distinct elements $i$ and $j$ of $[n]$, the
transposition $t_{i,j} \in S_n$ has sign $(-1)^{t_{i,j}} = -1$.

\textbf{(c)} For any positive integer $k$ and any distinct elements
$i_1, i_2, \ldots, i_k \in [n]$, the $k$-cycle
$\operatorname{cyc}_{i_1, i_2, \ldots, i_k}$ has sign
$(-1)^{\operatorname{cyc}_{i_1, i_2, \ldots, i_k}} = (-1)^{k-1}$.

\textbf{(d)} We have $(-1)^{\sigma\tau} = (-1)^{\sigma} \cdot (-1)^{\tau}$
for any $\sigma \in S_n$ and $\tau \in S_n$.

\textbf{(e)} We have $(-1)^{\sigma_1 \sigma_2 \cdots \sigma_p}
= (-1)^{\sigma_1} (-1)^{\sigma_2} \cdots (-1)^{\sigma_p}$
for any $\sigma_1, \sigma_2, \ldots, \sigma_p \in S_n$.

\textbf{(f)} We have $(-1)^{\sigma^{-1}} = (-1)^{\sigma}$
for any $\sigma \in S_n$. (The left hand side here has to be
understood as $(-1)^{(\sigma^{-1})}$.)

\textbf{(g)} We have
\[
(-1)^{\sigma} = \prod_{1 \leq i < j \leq n}
\frac{\sigma(i) - \sigma(j)}{i - j}
\qquad \text{for each } \sigma \in S_n.
\]
(The product sign ``$\prod_{1 \leq i < j \leq n}$'' means a product over
all pairs $(i,j)$ of integers satisfying $1 \leq i < j \leq n$.
There are $\binom{n}{2}$ such pairs.)

\textbf{(h)} If $x_1, x_2, \ldots, x_n$ are any elements of some
commutative ring, and if $\sigma \in S_n$, then
\[
\prod_{1 \leq i < j \leq n}
\left( x_{\sigma(i)} - x_{\sigma(j)} \right)
= (-1)^{\sigma}
\prod_{1 \leq i < j \leq n}
\left( x_i - x_j \right).
\]
\end{proposition}

\begin{proof}
Most of this follows easily from what we have proved above, but here are references to
complete proofs:

\textbf{(a)} This is \cite[Proposition 5.15 \textbf{(a)}]{detnotes}, and
follows easily from $\ell(\operatorname{id}) = 0$.

\textbf{(d)} This is \cite[Proposition 5.15 \textbf{(c)}]{detnotes}, and
follows easily from Corollary~\ref{cor.perm.red.sigtau}~\textbf{(a)}.

\textbf{(b)} This is \cite[Exercise 5.10 \textbf{(b)}]{detnotes}, and follows
easily from the fact that a permutation has even length if and only if
it has an even number of even-length cycles.

\textbf{(c)} This is \cite[Exercise 5.17 \textbf{(d)}]{detnotes}, and follows
easily from the fact that a cycle of length $\ell$ has exactly $\ell - 1$
inversions, combined with Proposition
\ref{prop.perm.sign.props}~\textbf{(d)}.

\textbf{(e)} This is \cite[Proposition 5.28]{detnotes}, and follows by
induction from Proposition~\ref{prop.perm.sign.props}~\textbf{(d)}.

\textbf{(f)} This is \cite[Proposition 5.15 \textbf{(d)}]{detnotes}, and
follows easily from Proposition~\ref{prop.perm.sign.props}~\textbf{(d)} or
from Proposition~\ref{prop.perm.len.inv}.

\textbf{(h)} This is \cite[Exercise 5.13 \textbf{(a)}]{detnotes} (or, rather,
the straightforward generalization of \cite[Exercise 5.13 \textbf{(a)}]{detnotes}
to arbitrary commutative rings). Each
factor $x_{\sigma(i)} - x_{\sigma(j)}$ on the left
hand side appears also on the right hand side, albeit with a different sign if
$(i,j)$ is an inversion of $\sigma$. Thus, the products on both
sides agree up to a sign, which is precisely $(-1)^{\ell(\sigma)} = (-1)^{\sigma}$.

\textbf{(g)} This is \cite[Exercise 5.13 \textbf{(c)}]{detnotes}, and is a
particular case of Proposition~\ref{prop.perm.sign.props}~\textbf{(h)}.
\end{proof}

\subsection{The sign homomorphism}

\begin{corollary}
\label{cor.perm.sign.hom}
\lean{Equiv.Perm.sign_hom_mul}
\leantarget
\leanok
Let $n \in \mathbb{N}$. The map
\begin{align*}
S_n &\to \{1, -1\}, \\
\sigma &\mapsto (-1)^{\sigma}
\end{align*}
is a group homomorphism from the symmetric group $S_n$ to the order-$2$
group $\{1, -1\}$. (Of course, $\{1, -1\}$ is a group with respect to
multiplication.)
\end{corollary}

This map is known as the \emph{sign homomorphism}.

\begin{proof}
Proposition~\ref{prop.perm.sign.props}~\textbf{(d)} shows that this map
respects multiplication (i.e., sends products to products). However, if a
map between two groups respects multiplication, then it is automatically a
group homomorphism. Thus, Corollary~\ref{cor.perm.sign.hom} follows.
\end{proof}

\subsection{Even and odd permutations}

\begin{definition}
\label{def.perm.even-odd}
\lean{Equiv.Perm.IsEven, Equiv.Perm.IsOdd}
\leantarget
\leanok
Let $n \in \mathbb{N}$. A permutation $\sigma \in S_n$ is said to be
\begin{itemize}
\item \emph{even} if $(-1)^{\sigma} = 1$ (that is, if $\ell(\sigma)$ is even);
\item \emph{odd} if $(-1)^{\sigma} = -1$ (that is, if $\ell(\sigma)$ is odd).
\end{itemize}
\end{definition}

\begin{lemma}
\label{lem.perm.isEven_or_isOdd}
\lean{Equiv.Perm.isEven_or_isOdd}
\leanhelper
\leanok
Every permutation $\sigma$ is either even or odd.
\end{lemma}

\begin{proof}
\leanok
$(-1)^{\sigma} \in \{1, -1\}$.
\end{proof}

\begin{lemma}
\label{lem.perm.not_isEven_and_isOdd}
\lean{Equiv.Perm.not_isEven_and_isOdd}
\leanhelper
\leanok
A permutation cannot be both even and odd.
\end{lemma}

\begin{proof}
\leanok
$1 \neq -1$.
\end{proof}

\begin{lemma}
\label{lem.perm.isEven_one}
\lean{Equiv.Perm.isEven_one}
\leanhelper
\leanok
The identity permutation is even.
\end{lemma}

\begin{proof}
\leanok
By Proposition~\ref{prop.perm.sign.props}~\textbf{(a)},
$(-1)^{\operatorname{id}} = 1$.
\end{proof}

\begin{lemma}
\label{lem.perm.isOdd_swap}
\lean{Equiv.Perm.isOdd_swap}
\leanhelper
\leanok
A transposition of two distinct elements is an odd permutation.
\end{lemma}

\begin{proof}
\leanok
By Proposition~\ref{prop.perm.sign.props}~\textbf{(b)},
$(-1)^{t_{i,j}} = -1$.
\end{proof}

\begin{lemma}
\label{lem.perm.isEven_iff_mem_alternatingGroup}
\lean{Equiv.Perm.isEven_iff_mem_alternatingGroup}
\leanhelper
\leanok
A permutation $\sigma$ is even if and only if $\sigma \in A_n$.
\end{lemma}

\begin{proof}
\leanok
This follows directly from the definition of $A_n$ as the kernel of the sign homomorphism:
$\sigma \in A_n$ iff $\operatorname{sign}(\sigma) = 1$ iff $\sigma$ is even.
\end{proof}

\begin{lemma}
\label{lem.perm.isOdd_iff_not_mem_alternatingGroup}
\lean{Equiv.Perm.isOdd_iff_not_mem_alternatingGroup}
\leanhelper
\leanok
A permutation $\sigma$ is odd if and only if $\sigma \notin A_n$.
\end{lemma}

\begin{proof}
\leanok
Since every permutation is either even or odd
(Lemma~\ref{lem.perm.isEven_or_isOdd}), and $\sigma$ is even iff
$\sigma \in A_n$ (Lemma~\ref{lem.perm.isEven_iff_mem_alternatingGroup}),
it follows that $\sigma$ is odd iff $\sigma \notin A_n$.
\end{proof}

\subsection{The alternating group}

\begin{corollary}
\label{cor.perm.altgp}
\lean{Equiv.Perm.alternatingGroup_isNormal}
\leantarget
\leanok
Let $n \in \mathbb{N}$. The set of all even permutations in $S_n$ is a
normal subgroup of $S_n$.
\end{corollary}

This subgroup is known as the $n$-th \emph{alternating group} (commonly
called $A_n$).

\begin{proof}
The set of all even permutations in $S_n$ is the kernel of the group
homomorphism $S_n \to \{1, -1\}$ from Corollary~\ref{cor.perm.sign.hom}.
Thus, it is a normal subgroup of $S_n$ (since any kernel is a normal
subgroup).
\end{proof}

\subsection{Counting even and odd permutations}

\begin{corollary}
\label{cor.perm.num-even}
\lean{Equiv.Perm.card_even_perms, Equiv.Perm.card_odd_eq_card_even}
\leantarget
\leanok
Let $n \geq 2$. Then,
\[
\left(\text{\# of even permutations } \sigma \in S_n \right)
= \left(\text{\# of odd permutations } \sigma \in S_n \right)
= n!/2.
\]
\end{corollary}

\begin{proof}
The symmetric group $S_n$ contains the simple transposition $s_1$ (since
$n \geq 2$). If $\sigma \in S_n$, then
by Prop.~\ref{prop.perm.sign.props}~\textbf{(d)}, we get
\[
(-1)^{\sigma s_1}
= (-1)^{\sigma} \cdot \underbrace{(-1)^{s_1}}_{
\substack{= -1 \\ \text{(by Prop.~\textbf{(b)},} \\
\text{since } s_1 = t_{1,2}\text{)}}}
= -(-1)^{\sigma},
\]
where the underbrace uses Prop.~\ref{prop.perm.sign.props}~\textbf{(b)}.
Hence, a permutation $\sigma \in S_n$ is even if and only if $\sigma s_1$
is odd. Hence, the map
\begin{align*}
\{\text{even permutations } \sigma \in S_n\}
&\to \{\text{odd permutations } \sigma \in S_n\}, \\
\sigma &\mapsto \sigma s_1
\end{align*}
is well-defined. This map is furthermore a bijection (since $S_n$ is a
group). Thus, the bijection principle yields
\[
\left(\text{\# of even permutations } \sigma \in S_n \right)
= \left(\text{\# of odd permutations } \sigma \in S_n \right).
\]
Both sides of this equality must furthermore equal $n!/2$, since they add
up to $|S_n| = n!$. This proves Corollary~\ref{cor.perm.num-even}.
(See \cite[Exercise 5.4]{detnotes} for details.)
\end{proof}

As a consequence of Corollary~\ref{cor.perm.num-even}, we see that
\begin{equation}
\label{eq.cor.perm.num-even.sum-sign}
\sum_{\sigma \in S_n} (-1)^{\sigma} = 0
\qquad \text{for each } n \geq 2.
\end{equation}
(Indeed, the sum $\sum_{\sigma \in S_n} (-1)^{\sigma}$ can be
rewritten as
\[
\left(\text{\# of even permutations } \sigma \in S_n \right)
- \left(\text{\# of odd permutations } \sigma \in S_n \right),
\]
since the addends corresponding to the even permutations $\sigma \in S_n$ are
equal to $1$ whereas the addends corresponding to the odd permutations
$\sigma \in S_n$ are equal to $-1$.)

\begin{lemma}
\label{lem.perm.sum-sign-eq-zero}
\lean{Equiv.Perm.sum_sign_eq_zero}
\leanhelper
\leanok
For $n \geq 2$,
$\displaystyle\sum_{\sigma \in S_n} (-1)^{\sigma} = 0$.
\end{lemma}

\begin{proof}
\leanok
This is a restatement of \eqref{eq.cor.perm.num-even.sum-sign}.
\end{proof}

\subsection{Sign for permutations of arbitrary finite sets}

We note that the sign can be defined not only for a permutation
$\sigma \in S_n$, but also for any permutation of any finite set $X$
(even if the set $X$ has no chosen total order on it, as the set $[n]$ has).
Here is one way to do so:

\begin{proposition}
\label{prop.perm.sign.X}
\lean{Equiv.Perm.sign_conj_eq, Equiv.Perm.sign_id_finiteSet, Equiv.Perm.sign_mul_finiteSet}
\leantarget
\leanok
Let $X$ be a finite set. We want to define the sign of any permutation of $X$.

Fix a bijection $\phi : X \to [n]$ for some $n \in \mathbb{N}$. (Such a
bijection always exists, since $X$ is finite.) For every permutation
$\sigma$ of $X$, set
\[
(-1)_{\phi}^{\sigma} := (-1)^{\phi \circ \sigma \circ \phi^{-1}}.
\]
Here, the right hand side is well-defined, since
$\phi \circ \sigma \circ \phi^{-1}$ is a permutation of $[n]$. Now:

\textbf{(a)} This number $(-1)_{\phi}^{\sigma}$ depends only on the
permutation $\sigma$, but not on the bijection $\phi$. (In other words,
if $\phi_1$ and $\phi_2$ are two bijections from $X$ to $[n]$, then
$(-1)_{\phi_1}^{\sigma} = (-1)_{\phi_2}^{\sigma}$.)

Thus, we shall denote $(-1)_{\phi}^{\sigma}$ by $(-1)^{\sigma}$ from
now on. We refer to this number $(-1)^{\sigma}$ as the \emph{sign} of the
permutation $\sigma \in S_X$. (When $X = [n]$, this notation does not
clash with Definition~\ref{def.perm.sign}, since we can pick the bijection
$\phi = \operatorname{id}$ and obtain
$(-1)_{\phi}^{\sigma} = (-1)^{\operatorname{id} \circ \sigma \circ
\operatorname{id}^{-1}} = (-1)^{\sigma}$.)

\textbf{(b)} The identity permutation $\operatorname{id} : X \to X$
satisfies $(-1)^{\operatorname{id}} = 1$.

\textbf{(c)} We have $(-1)^{\sigma\tau} = (-1)^{\sigma} \cdot (-1)^{\tau}$
for any two permutations $\sigma$ and $\tau$ of $X$.
\end{proposition}

\begin{proof}
This all follows quite easily from
Proposition~\ref{prop.perm.sign.props}~\textbf{(d)}. See
\cite[Exercise 5.12]{detnotes} for a detailed proof.
\end{proof}
